% !TEX root = main.tex

\begin{frame}
    \frametitle{Conclusions \& Future Work}

    % Testo precedente:
    % Conclusions
    % - Contract-driven assessment verifies security behavior beyond syntactic conformance.
    % - Specified oracles and the eight-domain taxonomy organize semantic security checks.
    % - The deterministic engine distinguishes enforced from violated guarantees and supports reproducible continuous verification.
    % Future Work
    % - Complete the coverage of the taxonomy and add adapters for managed cloud gateways.
    % - Develop reliable methods to generate security oracles automatically.
    % Solo nello speech: i limiti (oracoli manuali, D5 senza test, un solo target), ognuno legato a un future work.
    Contract-driven security assurance for Big Data APIs, with \alert{specified oracles}
    \begin{itemize}
        \item Verifies security \alert{behavior}, not only syntax
        \item \alert{Same outcomes} across independent runs $\rightarrow$ ready for CI/CD
        \item Found a real \alert{contract--behavior discrepancy} on Forgejo
    \end{itemize}

    \vspace{.8em}
    \uncover<2->{
            \textbf{Future Work}
            \begin{itemize}
                \item \alert{Automatic generation} of security oracles
                \item Complete the coverage of the taxonomy
                \item Validation on \alert{other targets}
            \end{itemize}
        }

\end{frame}

\begin{frame}

    % Il blocco "Thanks! Questions?" usa i colori e il carattere del titolo (stesso aspetto della barra
    % del titolo del primo frame). Sotto, i contatti: sono gli stessi di main.tex (nome, email, sito) e si scrivono due volte.
    \begin{center}
        \begin{beamercolorbox}[sep=.75em]{title}
            \centering
            \usebeamerfont{title}Thanks! Questions?%
        \end{beamercolorbox}

        \vspace{2em}

        {\large Enea Manzi}\\
        \vspace{.25em}
        {\small
            \href{mailto:enea.manzi@unimi.it}{enea.manzi@unimi.it}\\
            \href{https://eneamanzi.github.io/}{https://eneamanzi.github.io/}\\
            \href{https://www.linkedin.com/in/eneamanzi}{https://www.linkedin.com/in/eneamanzi}
        }

    \end{center}

\end{frame}
